%!TEX root = ../template.tex
%%%%%%%%%%%%%%%%%%%%%%%%%%%%%%%%%%%%%%%%%%%%%%%%%%%%%%%%%%%%%%%%%%%%
%% chapter2.tex
%% NOVA thesis document file
%%
%% Chapter with the template manual
%%%%%%%%%%%%%%%%%%%%%%%%%%%%%%%%%%%%%%%%%%%%%%%%%%%%%%%%%%%%%%%%%%%%

\typeout{NT FILE chapter2.tex}%

\chapter{Literature review}
This chapter reviews relevant literature on Multi-Agent Deep Deterministic Policy Gradient (MADDPG) and Graph Neural Networks (GNNs),
We will explore their architectures, the proposed applications in network routing and current literature on their vulnerabilities to adversarial attacks.
\section{DRL definition and overview}

\textbf{Deep Reinforcement Learning (DRL)} is a subfield of machine learning that combines 
reinforcement learning (RL) principles with deep neural networks to enable 
agents to learn optimal policies in complex environments. In DRL, 
an agent interacts with an environment by taking actions based on its current state, 
receiving feedback in the form of rewards or penalties. The goal of the agent is to 
learn a policy that maximizes the cumulative reward over time. 
Unlike traditional rule-based approaches, DRL agents learn 
directly from state observations, making them particularly 
suited for network optimization tasks where hand-crafted rules cannot 
capture the full complexity of system dynamics. \cite{theodoropoulos2023multi} \cite{10535707} \cite{lozanocuadra2024opensourcemultiagentdeep} 
In the context of network 
routing, DRL enables agents to discover adaptive forwarding policies that 
respond to real-time network conditions (congestion, link failures, 
traffic demands) without requiring explicit model of the environment. \cite{11227800} \cite{10730787} \cite{11122503}
The environment is naturally modeled as a Markov Decision Process 
(MDP) or Partially Observable MDP (POMDP) where states encode link loads and 
topology information, actions correspond to next-hop or path selection decisions,
and rewards capture delay, throughput, or composite QoS metrics. \cite{10684240}

However, DRL models are vulnerable to adversarial attacks, where malicious actors can
input adversarial examples designed to mislead the model into making incorrect decisions 
\cite{szegedy2014intriguingpropertiesneuralnetworks},these are often called either evasion attacks when introduced during test-time
or poisoning attacks when introduced during training. Lin et al. \cite{lin2019tacticsadversarialattackdeep} 
further introduce ‘enchanting attacks’, which craft a sequence of adversarial 
perturbations to drive the agent toward specific target states.
\cite{yamabe2025robustdeepreinforcementlearning} \cite{dai2025effectivebackdoorattackgraph}
These attacks exploit the inherent lack of robustness in the model, leading to suboptimal decision-making in many application areas. 
\cite{alzubaidi2025adversarial} \cite{zhou2024adversarial} \cite{finlayson2019adversarialexample}
Common adversarial attack methods include the Fast Gradient Sign Method (FGSM),
Projected Gradient Descent (PGD),the Carlini and Wagner (CW) attack and more recently,
the distribution-Aware Projected Gradient Descent (DAPGD). \cite{villegas2024evaluating} \cite{duan2025DADPG}


\section{Deep Q-Networks}
\textbf{Deep Q-Networks} integrate deep convolutional neural networks 
with Q-learning, a reinforcement learning algorithm, to learn 
optimal action-selection policies directly from high-dimensional sensory 
inputs without manual feature engineering. \cite{mnih2013playingatarideepreinforcement} \cite{mnih2015human} The architecture parameterizes 
the action-value function, the Q-function, using a deep neural network whose 
input consists of preprocessed high-dimensional observations 
(e.g., raw pixel data) and whose output comprises value estimates for each 
possible action, representing the expected cumulative future reward associated
with each action selection. Training stability, historically a critical challenge 
when combining neural networks with temporal-difference learning, is achieved
through two key mechanisms: experience replay, a technique
wherein the network is trained on minibatches of randomly sampled 
experiences drawn from a memory buffer, thereby breaking temporal correlations 
in the observation sequence and smoothing shifts in the data distribution. \cite{mnih2013playingatarideepreinforcement} \cite{mnih2015human}
The network is trained via
stochastic gradient descent to minimize the temporal-difference error, 
defined by the squared Bellman residual between the predicted Q-value 
and the target value computed as the immediate reward plus the discounted 
maximum Q-value ofthe subsequent state as estimated by the target network. \cite{mnih2013playingatarideepreinforcement}

\begin{equation}
L_i(\theta_i) = \mathbb{E}_{s,a \sim \rho(\cdot)} \left[ (y_i - Q(s, a; \theta_i))^2 \right]
\end{equation}

where the target value is:

\begin{equation}
y_i = \mathbb{E}_{s' \sim \mathcal{E}} \left[ r + \gamma \max_{a'} Q(s', a'; \theta_{i-1}) \Big| s, a \right]
\end{equation}

\begin{itemize}

\item \textbf{$L_i(\theta_i)$}: Loss function at iteration $i$ with network parameters $\theta_i$

\item \textbf{$\theta_i$}: Network parameters (weights and biases) at iteration $i$

\item \textbf{$\theta_{i-1}$}: Network parameters from the previous iteration $i-1$ (held fixed during optimization of $L_i$)

\item \textbf{$\mathbb{E}_{s,a \sim \rho(\cdot)}$}: Expectation over state-action pairs $(s,a)$ sampled from the behaviour distribution $\rho(s,a)$

\item \textbf{$y_i$}: Target value (Bellman target) for iteration $i$

\item \textbf{$Q(s, a; \theta_i)$}: Predicted action-value function (Q-function) that estimates the expected cumulative discounted reward for taking action $a$ in state $s$ using network parameters $\theta_i$

\item \textbf{$r$}: Immediate scalar reward received from the environment after taking action $a$ in state $s$

\item \textbf{$\gamma$}: Discount factor; a scalar in $[0, 1]$ that weights the importance of future rewards relative to immediate rewards

\item \textbf{$\max_{a'} Q(s', a'; \theta_{i-1})$}: Maximum Q-value across all possible actions $a'$ in the successor state $s'$, computed using the previous iteration's parameters $\theta_{i-1}$

\item \textbf{$s$}: Current state, represented as a sequence of preprocessed observations (e.g., stacked pixel frames)

\item \textbf{$a$}: Action taken in state $s$; chosen from the discrete action set

\item \textbf{$s'$}: Successor state resulting from taking action $a$ in state $s$

\item \textbf{$\mathcal{E}$}: The environment (in the DQN paper: the Atari 2600 emulator \cite{mnih2013playingatarideepreinforcement}); transition dynamics are unknown and learned implicitly

\item \textbf{$\rho(s,a)$}: Behaviour distribution; the probability distribution of state-action pairs encountered during training (typically from an $\epsilon$-greedy policy)

\end{itemize}
\cite{mnih2013playingatarideepreinforcement}

DQN has been successfully applied to network routing tasks, 
where the agent learns to make optimal forwarding decisions based on various
real-time network state observations. A recent paper by Xuan Geng et al. (2025) \cite{10012647} Proposes a DQN-based
routing protocol that optimizes both path selection and energy efficiency thereby
significantly extending network lifetime of underwater acoustic sensor networks (UASNs).
DQL has also been proposed as a solution to help Software Defined Networks (SDN) based
data-centers dealing with "mice" flows, which are short-lived flows with realtively small
data sizes that can congest the network if not handled properly and "elephant" flows,
that are long lived and carry large amounts of data. 2 DQNs are built, one for each type
of flow, that train Convulutional Neural Networks (CNN) to optimize routing decisions. \cite{9105030}

\subsection{Duelling Q-networks}
Dueling Network Architectures for Deep Reinforcement Learning represent a 
structural innovation to the standard Deep Q-Network architecture in which the 
action-value function is decomposed into two independent,
streams that share a common feature representation before final aggregation 
\cite{wang2016duelingnetworkarchitecturesdeep,zhao2025breakingcoldstartbarrierreinforcement}. Specifically, the architecture bifurcates into a value 
stream $V(s; \theta_v)$ that estimates the scalar value of a given state 
independent of action selection, and an advantage stream $A(s, a; \theta_a)$ that
computes the relative advantage of executing a particular action within that
state, with both streams processing shared convolutional features derived from
the input observation \cite{wang2016duelingnetworkarchitecturesdeep}. The final action-value estimate is
obtained through a dueling layer that performs a learned aggregation of these
two quantities and provides 
a normalization mechanism that stabilizes gradient-based learning and enforces 
identifiability \cite{wang2016duelingnetworkarchitecturesdeep}. This architectural separation facilitates
more efficient learning in domains where numerous actions yield 
comparatively similar expected returns, by enabling the value and advantage streams 
to specialize in capturing specific aspects of the policy structure while 
maintaining a shared feature space, thereby improving sample efficiency 
and convergence properties 
\cite{wang2016duelingnetworkarchitecturesdeep}.

In a recent paper by Zhang et al. (2022) \cite{ZHANG2022108663} a deep dueling Q-network based routing protocol
helps optimize Routing and Spectrum Allocation (RSA) in Elastic Optical Networks (EONs).
The proposed method leverages the dueling architecture to effectively learn 
efficient resource allocation policies under dynamic traffic 
conditions, where connection requests arrive and terminate unpredictably. The 
decomposition of the Q-function into value and advantage streams allows the model to
better judge the relative merits of different routing actions based on the 
overall network congestion and fragmentation state
thereby circumventing the learning inefficiency that arises when 
traditional DQN attempts to differentiate between numerous 
similar-valued routing actions across diverse network configurations.\cite{ZHANG2022108663}

\section{MADDPG}

\textbf{Multi-Agent Deep Deterministic Policy Gradient (MADDPG)} 
is a deep reinforcement learning algorithm designed for multi-agent environments 
that extends actor-critic methods through a centralized training with decentralized 
execution framework, where each agent maintains a centralized action-value function 
conditioned on agent actions during training while 
using only local observations at execution time. To relax the assumption of knowing other 
agent's policies, MADDPG enables each agent to maintain an approximate policy 
of other agents by maximizing the log-probability of observed actions 
with entropy regularization, 
allowing subject agents to learn target agents policies before adapting accordingly. For enhanced robustness in competitive settings, MADDPG trains a 
collection of different sub-policies per agent, randomly selecting one sub-policy 
per episode and maintaining separate replay buffers for each 
sub-policy , which maximizes the ensemble objective 
such that agents develop diverse behavioral strategies that prevent overfitting to specific 
opponent policies. \cite{lowe2017multi}

However, this very architecture creates exploitable attack surfaces. 
The centralized critic's dependence on gradient computations makes 
MADDPG amenable to gradient-based adversarial perturbations, while 
the decentralized execution phase means each agent operates 
independently at test time, creating coordination vulnerabilities 
when one or more agents receive adversarial inputs. \cite{10684240,pham2022c,li2025attacking}

\subsection{MADDPG in Network Routing}
A recent paper by Tian et al. (2023) \cite{Tian2023maddpgslice} proposes a MADDPG-based
approach for solving the network slicing problem in 5G networks, a resource
allocation problem where diverse service types with heterogeneous QoS requirements
may, at any point, change their resource demands. The proposed method, MADDPG-SR
,treats each virtual network function (VNF) as an agent that learns to allocate resources
and migrate dynamically with the end goal of maintaining QOS requirements in every slice.
The results show the proposed method outperforms non-MADDPG implementations. \cite{Tian2023maddpgslice}

A recent paper by Huang et al. (2025) \cite{huang2025era} proposes ERA-MADDPG as an elastic routing scheme for 
software-defined networks (SDNs) that casts routing optimization
as a cooperative multi-agent decision-making problem and uses 
Multi-Agent Deep Deterministic Policy Gradient (MADDPG) as its 
learning backbone. In this framework, network entities 
(e.g., switches or abstract routing agents) learn deterministic 
policies that select paths or adjust forwarding behavior based on 
observed traffic and topology information, while a centralized critic
exploits global state during training to stabilize learning in the
inherently non-stationary multi-agent environment. The resulting 
actors are then deployed in a decentralized fashion, allowing each 
agent to make fast local decisions that collectively adapt the 
end-to-end routes to time-varying traffic conditions and QoS demands.
By leveraging MADDPG, ERA-MADDPG overcomes key limitations of static 
routing and single-agent RL methods, notably poor scalability and weak 
coordination, enabling continuous “elastic” adjustment of routes to 
jointly reduce delay, packet loss, and congestion while improving load
balance. This method outperforms other DRL based algorithms on training
efficiency and recovery in dynamic network topologies. \cite{huang2025era}

Another recent design, RS-MADDPG, \cite{kuang2025rs} enables each node to learn a deterministic policy 
that maps local observations to continuous routing decisions 
(e.g., flow splitting ratios or next-hop selections) while 
implicitly accounting for the actions of other agents, thereby 
mitigating the non-stationarity that plagues independent RL approaches.
By leveraging MADDPG's capacity for continuous action spaces
and multi-agent credit assignment, RS-MADDPG dynamically adjusts 
routes in response to real-time traffic patterns, link failures, 
and varying QoS demands without requiring manual reconfiguration. 
The algorithm's effectiveness is further enhanced by its ability to 
converge approximately 30 training cycles earlier than baseline 
methods, demonstrating that the Centralized Training Decentralized execution
(CTDE) paradigm, caracteristic of MADDPG, accelerates learning 
of cooperative policies that jointly minimize latency, balance loads,
and maximize throughput. Consequently, RS-MADDPG not only 
outperforms conventional protocols in key performance metrics
but also provides a scalable, adaptive solution for delivering
differentiated QoS services in modern SDN environments. \cite{kuang2025rs}

An interesting development based on the MADDPG algorithm is the 
Dynamic Packet Routing Algorithm using Multi-Agent Reinforcement Learning with Emergent Communication (DRAMA)
\cite{zhang2025dramadynamicpacketrouting}. DRAMA utlizes the Centralized
Training Decentralized Execution (CTDE) framework of MADDPG,
but redesign the value-function architecture and information
flow to overcome MADDPG's key limitations in packet routing.
While MADDPG uses a Type‑II critic that takes explicit joint
observation, action pairs as input, requiring fixed, enumerated
discrete actions and becoming brittle when the action set 
changes, DRAMA replaces this with a Q-network score layer that
operates on topology-aware latent features extracted by an
observation encoder and a graph-based emergent communication
layer. This preserves the CTDE advantage
of learning from joint observations and rewards while decoupling
the Q-function from a fixed action index space, so the same trained
network can seamlessly score newly appearing neighbors or cope with
link/router failures without the need for extensive retraining. Empirically,
this MADDPG-based CTDE training combined with DRAMA's communication
and Q-network design yields comparatively better delivery rates and
lower end-to-end latency than regular MADDPG implementations across
a range of traffic loads and under dynamic topologies,
demonstrating that DRAMA leverages the strengths of MADDPG's training
paradigm while addressing its natural difficulty with
scalable, adaptive routing. \cite{zhang2025dramadynamicpacketrouting} 

Although the usage of MADDPG in network routing has increased in recent years,
there is lack of research on the effects of adversarial attacks in this context.

\section{Taxonomy on Adversarial attacks on MADDPG}

Adversarial attacks against MADDPG can be systematically categorized
along three dimensions: (1) perturbation domain 
(state, action, communication, reward, model), (2) attack timing 
(training-time vs. test-time), and (3) attacker knowledge 
(white-box, gray-box, black-box).

\begin{tikzpicture}[
    node distance=0.6cm and 0.8cm,
    every node/.style={font=\small},
    box/.style={rectangle, draw, rounded corners, minimum width=1.4cm, minimum height=0.5cm, align=center},
    root/.style={box, fill=blue!20, font=\bfseries\tiny},
    category/.style={box, fill=cyan!15, font=\tiny},
    attack/.style={box, fill=green!10, font=\tiny, minimum width=1.2cm},
    ]
    
    % Root node
    \node[root] (root) {Adversarial\\Attacks on\\MADDPG};
    
    % Category 1: Observation Space
    \node[category, below=of root, xshift=-4.5cm, yshift=-1cm] (obs) {Observation\\Space};
    \draw[-latex] (root) -- (obs);
    \node[attack, below=of obs] (fgsm) {FGSM/\\PGD};
    \node[attack, below=of fgsm] (saja) {SAJA};
    \node[attack, below=of saja] (paad) {PA-AD};
    \draw[-latex] (obs) -- (fgsm);
    \draw[-latex] (obs) -- (saja);
    \draw[-latex] (obs) -- (paad);
    
    % Category 2: Action Space
    \node[category, below=of root, xshift=-2.7cm, yshift=-1cm] (action) {Action\\Space};
    \draw[-latex] (root) -- (action);
    \node[attack, below=of action] (direct) {Direct\\Action};
    \node[attack, below=of direct] (wolf) {Wolfpack};
    \node[attack, below=of wolf] (ami) {AMI};
    \draw[-latex] (action) -- (direct);
    \draw[-latex] (action) -- (wolf);
    \draw[-latex] (action) -- (ami);
    
    % Category 3: Communication Space
    \node[category, below=of root, yshift=-1cm] (comm) {Comm.\\Space};
    \draw[-latex] (root) -- (comm);
    \node[attack, below=of comm] (msg) {Message\\Hacking};
    \node[attack, below=of msg] (adacomm) {Adaptive\\Comm.};
    \draw[-latex] (comm) -- (msg);
    \draw[-latex] (comm) -- (adacomm);
    
    % Category 4: Poisoning & Backdoor
    \node[category, below=of root, xshift=2.7cm, yshift=-1cm] (poison) {Poisoning\\Backdoor};
    \draw[-latex] (root) -- (poison);
    \node[attack, below=of poison] (modpoi) {Model\\Poisoning};
    \node[attack, below=of modpoi] (bdoor) {Backdoor};
    \draw[-latex] (poison) -- (modpoi);
    \draw[-latex] (poison) -- (bdoor);
    
    % Category 5: Transferability & Black-Box
    \node[category, below=of root, xshift=4.5cm, yshift=-1cm] (transfer) {Transfer\\Black-Box};
    \draw[-latex] (root) -- (transfer);
    \node[attack, below=of transfer] (domain) {Domain\\Random.};
    \node[attack, below=of domain] (bami) {Black-Box\\AMI};
    \draw[-latex] (transfer) -- (domain);
    \draw[-latex] (transfer) -- (bami);
    
\end{tikzpicture}


\subsection{Observation Space Attacks}
State-space attacks perturb the agent's observations
before the actor network processes them. These attacks exploit the
gradient-based learning mechanisms underlying neural network policies.

\textbf{Gradient-Based State Perturbations}: The foundational approach 
applies the Fast Gradient Sign Method (FGSM) or Projected Gradient 
Descent (PGD) to compute adversarial observations. 
FGSM performs a single-step perturbation in the gradient direction 
of the loss function,\cite{goodfellow2015explainingharnessingadversarialexamples} while PGD extends this through iterative 
optimization with projection \cite{madry2019deeplearningmodelsresistant}. However, according to Guo et al.\cite{guo2025sajastateactionjointattack} standard FGSM attacks on 
MADDPG agents show limited effectiveness because the critic's 
Q-value gradients alone provide unreliable guidance, imperfect critic 
training leads to noisy gradients. \cite{Waghela_2024}

\textbf{State-Action Joint Attacks (SAJA)}: Recent advances 
recognize that state and action perturbations interact 
synergistically. The SAJA framework operates in two phases: 
(1) a multi-step gradient ascent on actor and critic networks to 
compute adversarial states, then (2) a second gradient ascent using 
the perturbed state to craft adversarial actions. 
SAJA employs a balanced loss function combining negative Q-value 
with an action distance term (Maximum Action Difference—MAD), 
which proves that maximizing the multi-agent 
action difference is a principled objective for degrading team value. \cite{guo2025sajastateactionjointattack}

Empirical results demonstrate SAJA's effectiveness: reward drops 
reach up to 9.90\% on 3-agent FACMAC, and SAJA bypasses defenses 
including M3DDPG, PAAD, and ATLA. The synergy ratio of 
joint attack impact to sum of individual attack impacts exceeds one
confirming that joint state-action perturbations are more potent. \cite{guo2025sajastateactionjointattack}

\textbf{Policy Adversarial Actor-Director (PA-AD)}: Rather than directly 
learning state perturbations in high-dimensional spaces, 
PA-AD proposes a Policy Adversary MDP (PAMDP) where an optimal 
policy induces the optimal state adversary. The attack employs 
collaborative "director" and "actor" components: the director learns 
which policy perturbation directions are most effective, while the 
actor crafts state perturbations to realize those directions. 
PA-AD achieves state-of-the-art performance on Atari and MuJoCo 
tasks, and combined with adversarial training, produces robust 
policies that significantly outperform baseline defenses.\cite{sun2021strongest}

\subsection{Action-Space Attacks}
Action-space attacks directly manipulate an agent's output actions, 
either through direct perturbation of action values or by injecting 
adversarial agents that execute perturbed actions.

\textbf{Direct Action Manipulation}: SAJA's action-only variant applies 
gradient ascent on the critic network conditioned on the perturbed 
state to craft final adversarial actions. These attacks are 
particularly effective when attackers possess partial control over
one or more agent's actuators. \cite{guo2025sajastateactionjointattack}

\textbf{Malicious Agent Injection}: Rather than perturbing a single agent's 
actions, attackers can inject fully adversarial agents into the 
multi-agent system. These agents learn to disrupt the victim agent's 
policies through direct behavioral interference. The Wolfpack attack 
exemplifies this: by targeting a single agent and subsequently 
attacking its supporting agents, adversaries exploit the inherent 
interdependencies in cooperative systems. The Wolfpack approach 
reduces rewards by up to 51.8\% and 
severely impacts SMAC scenarios, proving far more effective than 
random or single-agent attacks. \cite{lee2025wolfpackadversarialattackrobust}

\textbf{Adversarial Minority Influence (AMI)}: AMI leverages social 
psychology principles where a minority (adversary) unilaterally 
sways the majority (victims) toward worst-case failure. AMI employs 
two innovations: (1) a unilateral influence filter adapted from 
mutual information that maximizes adversary-to-victim influence 
while excluding reverse influence, and (2) a targeted adversarial 
oracle that learns worst-case target actions. AMI successfully 
attacks real robot swarms and simulation environments including 
StarCraft II and multi-agent MuJoCo. \cite{li2025attacking}

\subsection{Communication-Space Attacks}
Multi-agent systems with explicit communication channels face a 
distinct attack surface: adversarial manipulation of inter-agent 
messages.

\textbf{Message Hijacking}: The paper "Hijacking Robot Teams Through 
Adversarial Communication" by Wu et al.\cite{wu2023hijacking} demonstrates that an adversarial policy 
can intercept and modify messages between agents. The authors report several‑fold 
performance degradation (e.g., more than four‑fold in some simulated predator–prey 
tasks and roughly a factor of two in Robotarium experiments), illustrating that 
adversarial communication can catastrophically disrupt team performance.
The attack bypasses agents' native policies because the manipulated 
messages are indistinguishable from benign communication within the 
agents' learned policies.

\textbf{Communication-Robust Learning}: Advanced attacks model adversarial 
communication as a cooperative problem, learning adaptive attack 
strategies that persistently disrupt coordination. The work on 
differentially private communication reveals that privacy-preserving 
mechanisms, while protecting against eavesdropping, do not 
inherently defend against adversarial message perturbations.\cite{zhao2023dpmac,yuan2023communicationrobustmultiagentlearningadaptable}

\subsection{Poisoning and Backdoor Attacks}
Training-time attacks embed long-term vulnerabilities into learned 
policies. \cite{dai2025effectivebackdoorattackgraph}

\textbf{Model Poisoning via MARL Coordination}: In federated learning with 
multiple poisoning attackers, attackers coordinate using MADDPG 
itself to craft effective global model updates that degrade overall 
system performance. The attack assumes attackers can approximate the
benign data distribution, then formulate the poisoning problem as an 
MDP where attackers jointly learn attack policies through MADDPG, 
outperforming baseline poisoning methods. \cite{shen2021coordinated}

\textbf{Backdoor Attacks (TrojDRL, BACKDOORL, BadRL, MARLNet)}: These 
frameworks embed trigger conditions during training. Once triggered, 
agents execute attacker-specified actions while appearing normal 
otherwise. MARLNet specifically targets multi-agent scenarios using 
imperceptible triggers and reward manipulation. Advanced backdoor 
frameworks achieve near-perfect attack success rates (ASR approximating 99\%) 
with minimal utility loss on benign tasks. Backdoor Vaccination, 
where defenders inject controlled triggers—offers a proactive 
defense by leveraging the backdoor concept itself. \cite{dai2025effectivebackdoorattackgraph,guo2025pnact,guo2023policycleanse,chen2022backdoorattacksmultiagentcollaborative}

\subsection{Transferability and Black-Box Attacks}

\textbf{Domain Randomization for Transfer}: Adversarial attacks trained on 
diverse source environments transfer to unseen target policies. 
Domain randomization across action spaces, reward functions, 
transition dynamics, policy algorithms, and network structures 
enables black-box or gray-box attacks without knowledge of target 
specifics. Interestingly, attacks trained on randomized domains 
sometimes outperform white-box attacks on the target, suggesting 
that robust attacks develop from exposure to environmental diversity.
\cite{pan2020transferable}

Adversarial Minority Influence (Black-Box): AMI operates without 
model access, learning attack policies through repeated interaction. 
The unilateral influence filter enables the attacker to focus on 
influencing victims while disregarding reverse effects, making the 
attack practical in black-box settings.\cite{li2025attacking}



\subsection{FGSM}
\textbf{Fast Gradient Sign Method (FGSM)} \cite{goodfellow2015explainingharnessingadversarialexamples} is an efficient single-step gradient based
technique for generating 
adversarial examples, it leverages the adaptive nature of neural networks to create small perturbations
that maximize the loss increase. It is defined as follows:
\begin{equation}
\eta = \epsilon \cdot \text{sign}(\nabla_x J(\theta, x, y))
\end{equation}

\noindent where:
\begin{itemize}
    \item $\theta$ denotes the model parameters
    \item $x$ represents the input
    \item $y$ is the target label
    \item $J(\theta, x, y)$ is the cost function used to train de model
\end{itemize}

The resulting adversarial example is obtained as:

\begin{equation}
\tilde{x} = x + \epsilon \cdot \text{sign}(\nabla_x J(\theta, x, y))
\end{equation}

\noindent where $\epsilon$ is a small scalar controlling perturbation magnitude, 
chosen to remain imperceptible to human observers while remaining sufficient to fool 
the classifier. 

This means that FGSM not only misleads models but makes them confident
in the mistakes that they make, successfully "convincing" them that their
missclassifications are correct with a significant success rate\cite{villegas2024evaluating}. Goodfellow et al. 
\cite{goodfellow2015explainingharnessingadversarialexamples} propose 
that FGSM's efficacy arises from the inherent linearity of neural 
networks in high-dimensional spaces, fundamentally contradicting 
the prevailing hypothesis that adversarial vulnerabilities stem 
from deep network non-linearity.

FGSM efficiency and simplicity have made it a widely used baseline for
development of a wide range of sophisticated adversarial attack methods. \cite{akhtar2021advances}

Rozsa et al. (2016) \cite{rozsa2016adversarial} extends FGSM by introducing
the hot/cold (HC) approach, the hot/cold method operates at the 
penultimate layer, constructing a feature vector that simultaneously 
boosts value of features of a target "hot" class and reduces value of features of the original "cold" class
enabling backpropagation to generate perturbations that push samples 
toward different target classes.

Kurakin et al. (2017) \cite{kurakin2017adversarialmachinelearningscale} proposes the Iterative
Fast Gradient Sign Method (I-FGSM), the core methodology generates 
adversarial examples dynamically during training: at each iteration, 
adversarial examples are produced from clean samples using the current 
network state and mixed into minibatches with clean examples,
enabling balanced gradient updates through a weighted loss function.

Xie et al. (2019) \cite{Xie_2019_CVPR} proposes Diverse Inputs Iterative 
Fast Gradient Sign Method (DI²-FGSM), a principled extension of 
FGSM that addresses the white-box-black-box transferability trade-off 
inherent in the FGSM family of attacks. Rather than directly 
maximizing the loss with respect to original inputs, as in I-FGSM \cite{kurakin2017adversarialmachinelearningscale}, 
DI²-FGSM applies stochastic, differentiable transformations 
(random resizing and padding) at each iteration with probability p
, maximizing loss with respect to these transformed inputs.


\subsection{PGD}
\textbf{Projected Gradient Descent (PGD)} is an iterative adversarial attack method that 
extends FGSM by applying multiple smaller perturbation steps.
\noindent
Projected Gradient Descent (PGD) is introduced by Madry et al. \cite{madry2019deeplearningmodelsresistant} 
as a multi-step, adversarial attack that iteratively updates an input within a constrained perturbation set to 
maximize loss. Concretely, for an allowed perturbation 
set $S$ (typically an $\ell_\infty$-ball around a clean input $x$), PGD starts
from an initial point $x_0 \in x + S$ and repeatedly applies gradient ascent on 
the loss with respect to the input followed by projection back onto $x+S$ following
the negative loss function:
\[
x_{t+1} = \Pi_{x+S}\bigl(x_t + \alpha\,\mathrm{sign}(\nabla_x L(\theta, x_t, y))\bigr),
\]
thereby seeking an adversarial example that attains high loss while 
respecting the perturbation budget.
They define budget types in the following categories:
\begin{itemize}
    \item $L_\infty$: This constraint sets a maximum modification budget for every pixel. By this definition FGSM is an $L_\infty$ single-step version of PGD.
    \item $L_2$: This constraint sets a maximum modification budget based on the Euclidean distance between the original and adversarial images.
    \item $L_1$: This constraint sets a maximum modification budget based on a set limit of which the sum of all modifications cannot exceed. (Not tested but implicitly compatible with PGD)
\end{itemize}

Through extensive experiments with random restarts, the authors demonstrate that PGD reliably finds local 
maxima whose loss values are tightly concentrated, and thus regard PGD as 
a "universal first-order adversary", meaning that 
robustness to PGD is indicative of robustness to any attack that only uses 
first-order (gradient) information. \cite{madry2019deeplearningmodelsresistant}

A survey by Akhtar et al. (2021) \cite{akhtar2021advances} highlights a fenomenon called
"label leaking" observed in FGSM adversarial training where the model actually 
overfits for adversarial images resulting in higher accuracy on adversarial examples
than on clean images. PGD's method avoids this issue by using random starts that increase
variety of modifications. \cite{madry2019deeplearningmodelsresistant}

\newpage
\subsection{Overview of adversarial attacks}
This subsection is for a summary and ease of reference of the attacks 
mentioned previously. This is by no means an exhaustive list of adversarial
attacks, but should give a good overview of the state of the art.

\begin{table}[h!]
\centering
\caption{Taxonomy of Adversarial Attacks on MADDPG}
\label{tab:maddpg_attacks}
\begin{tabular}{|p{2.5cm}|p{3cm}|p{4cm}|p{3.5cm}|}
\hline
\textbf{Attack Domain} & \textbf{Attack Name} & \textbf{Mechanism} & \textbf{Efficacy} \\
\hline
\multirow{4}{2.5cm}{Observation Space} & FGSM/PGD & Single/iterative gradient-based state perturbations & Limited effectiveness; imperfect critic gradients \\
\cline{2-4}
 & SAJA & Two-phase gradient ascent on state then action with Maximum Action Difference (MAD) loss & 9.90\% reward drop (3a-FACMAC); synergy ratio $> 1$; bypasses M3DDPG, PAAD, ATLA \\
\cline{2-4}
 & PA-AD & Policy Adversary MDP with director-actor collaboration & SOTA on Atari/MuJoCo; outperforms FGSM/PGD \\
\hline
\multirow{3}{2.5cm}{Action Space} & Direct Action Manipulation & Gradient ascent on critic for action perturbations & Effective with partial actuator control; synergistic with state attacks \\
\cline{2-4}
 & Wolfpack Attack & Coordinated multi-stage targeting; exploit inter-agent dependencies & Up to 51.8\% reward reduction; severe SMAC impact \\
\cline{2-4}
 & AMI (Adversarial Minority Influence) & Unilateral influence filter + targeted adversarial oracle & Attacks robot swarms, StarCraft II, MuJoCo; black-box capable \\
\hline
\multirow{2}{2.5cm}{Communication Space} & Message Hijacking & Intercept and modify inter-agent messages & 465\% reduction (sim); 201\% (physical robots on Robotarium) \\
\cline{2-4}
 & Adaptive Communication Attacks & Learn persistent coordination disruption strategies & Bypasses privacy mechanisms; adaptive to learned patterns \\
\hline
\end{tabular}
\end{table}

\begin{table}[h!]
\centering
\caption{Taxonomy of Adversarial Attacks on MADDPG part 2}
\label{tab:maddpg_attacks2}
\begin{tabular}{|p{2.5cm}|p{3cm}|p{4cm}|p{3.5cm}|}
\hline
\textbf{Attack Domain} & \textbf{Attack Name} & \textbf{Mechanism} & \textbf{Efficacy} \\
\hline
\multirow{3}{2.5cm}{Poisoning \& Backdoor} & Model Poisoning (Federated) & Attackers coordinate via MADDPG to craft poisoned global updates & Outperforms baseline poisoning; coordinated attack via MDP \\
\cline{2-4}
 & Backdoor Attacks & Embed imperceptible triggers during training; execute on trigger & ASR $\approx 99\%$; minimal utility loss on benign tasks \\
\cline{2-4}
 & PNAct (Safe RL Backdoor) & Positive/negative action samples toggling safe/unsafe behavior & Stealthy; increases safety violations while maintaining reward \\
\hline
\multirow{2}{2.5cm}{Transfer \& Black-Box} & Domain Randomization Transfer & Train on diverse environments; generalize to unseen targets & Transfers across policies, networks, action spaces; sometimes outperforms white-box \\
\cline{2-4}
 & AMI (Black-Box) & Learn attack policies via repeated interaction without model access & No model knowledge required; effective on real swarms \\
\hline
\end{tabular}
\end{table}

\newpage

\section{Graph Neural Networks}
\textbf{Graph Neural Networks (GNNs)}
are a family of deep learning models designed to operate directly 
on graph-structured data by iteratively propagating and aggregating 
information along edges, so that each node (or graph) learns a 
representation that encodes both its own attributes and the topology 
of its neighborhood.Their core ideias are based on these principles:
Local connectivity, weight sharing, and hierarchical feature extraction. 
Modern surveys typically define GNNs as neural architectures that take a graph 
as input and learn task-specific node, edge, or graph-level outputs
via repeated “message passing” or neighborhood aggregation,
distinguishing them from shallow network embedding methods that merely
learn fixed vector lookups without parameter sharing across nodes. \cite{zhou2021graphneuralnetworksreview,zhang2022gnnreview,wu2021gnnreview,chami2022gnnreview}

Across the four surveys, the most common taxonomy divides GNNs into 
recurrent GNNs, convolutional GNNs, graph autoencoders, and 
spatio-temporal GNNs (with attention-based and message-passing 
variants largely falling under the convolutional category). 
Recurrent GNNs (often called “GNNs” in the early literature) 
define node states via recursive equations that iteratively update
each node’s hidden representation using its own previous state and 
the states of its neighbors until a fixed point is reached, followed 
by a task-specific readout function. 
Their main advantage is conceptual generality: 
they naturally handle arbitrary graph structures 
(directed/undirected, cyclic/acyclic) and can, in principle, 
approximate complex relational dependencies by iterating message 
passing to convergence. However, they are computationally expensive, 
require contraction conditions for guaranteed convergence, and are 
difficult to scale to large graphs because many recurrent steps and 
storage of intermediate states are needed; these limitations largely 
motivated the shift to feedforward, layer-stacked convolutional GNNs. 
\cite{zhou2021graphneuralnetworksreview,zhang2022gnnreview,wu2021gnnreview,chami2022gnnreview}

Convolutional GNNs (CGNNs) extend convolution from regular grids to graphs by 
defining each layer as a localized aggregation over a node’s neighborhood.
\cite{NIPS2015_536a76f9,tang2024chebnet} In the spectral view, this operation 
is grounded in graph signal processing, with explicit notions of frequency, 
smoothness, and graph filters. Polynomial filter approximations lead to spatially 
localized convolutions and practical implementations, but exact spectral methods 
are typically expensive and often restricted to transductive settings on a fixed 
graph.
\cite{zhou2021graphneuralnetworksreview,zhang2022gnnreview,wu2021gnnreview,chami2022gnnreview} As a result, classical spectral CGNNs offer strong theoretical foundations but are less natural to transfer across graphs with different topologies.

Spatial CGNNs (such as GraphSAGE and message‑passing neural networks) instead define 
convolutions directly in the node domain by aggregating neighbor features with 
permutation‑invariant functions (e.g., mean, sum, max, attention), often 
combined with neighborhood sampling to control complexity.\cite{jiawei2024graphsage++,raja2024artificial} Because the aggregation functions are shared across nodes and graphs, these models scale better to large graphs and naturally support inductive generalization to unseen nodes or entirely new graphs.\cite{zhou2021graphneuralnetworksreview,zhang2022gnnreview,wu2021gnnreview,chami2022gnnreview} At the same time, deep stacks of spatial layers can suffer from over‑smoothing, large‑graph performance can depend heavily on the sampling strategy, and the theoretical analysis is generally less mature than in strictly spectral formulations.

Graph autoencoders (GAEs) form another major family. 
An encoder GNN maps nodes or graphs into a latent space, 
and a decoder attempts to reconstruct structural information 
such as edges or the adjacency matrix, often under an unsupervised 
or self‑supervised loss.\cite{xu2024hc,gao2022gae} This view unifies 
many earlier network embedding methods within a single deep architecture: 
the learned representations capture topology and attributes and can then be 
reused for downstream tasks such as clustering, link prediction, or node 
classification.\cite{zhou2021graphneuralnetworksreview,zhang2022gnnreview,wu2021gnnreview,chami2022gnnreview} GAEs are particularly useful when labels are scarce, and variational extensions allow modeling distributions over graphs. Their main drawbacks are the cost of decoders that operate on dense adjacency matrices, the risk that pure reconstruction objectives over‑prioritize structural fidelity over task‑relevant semantics, and sensitivity to design choices such as negative sampling and the balance between observed and unobserved edges.

Spatio temporal GNNs (STGNNs) extend these architectures to graphs with time varying 
signals or topology by combining graph convolutions with temporal models 
such as recurrent networks, temporal convolutions, or attention mechanisms.
\cite{jin2022stgnn,yin2025stgnn} 
They are now a standard choice for applications like traffic 
forecasting and human action recognition, where the graph structure 
is relatively stable but node features evolve over time.\cite{wu2021gnnreview,zhou2021graphneuralnetworksreview} The main benefit is that STGNNs can model how information propagates through both space and time—for example, how congestion spreads across a road network—in an end‑to‑end trainable way. The cost is increased architectural and computational complexity, a greater need for labeled time‑series data, and more challenging regularization and interpretability due to the combined depth in the spatial and temporal dimensions.

Overall, recent surveys converge on the 
view that GNNs are message-passing architectures on graphs, 
with recurrent, convolutional, autoencoding, and spatio-temporal 
variants offering different trade-offs between expressivity, 
scalability, transferability, and ease of training. \cite{zhou2021graphneuralnetworksreview,zhang2022gnnreview,wu2021gnnreview,chami2022gnnreview}

\subsection{GNN applications in network routing}

Graph Neural Networks (GNNs) have emerged as a promising tool for learning routing and 
traffic engineering (TE) policies directly from network structure and traffic, with 
an explicit goal of generalizing across topologies and traffic conditions. Existing 
work can be organized into three main application axes: (i) GNNs as performance models 
for routing, (ii) GNN+RL controllers for online routing/TE, and (iii) specialized, 
generalizable GNN architectures for large-scale TE and domain-specific routing. \cite{jianggraph}

Early work focuses on using GNNs as surrogate models of network behavior, predicting 
performance metrics for given topologies, traffic variations, and routing configurations. 
Rusek et al. (2019) \cite{rusek2019unveiling} show that message-passing GNNs can 
accurately estimate per-link delay and loss, enabling fast evaluation of many 
candidate routing schemes without full simulation. Subsequent studies, summarized in 
Jiang et al. (2021) \cite{jianggraph} , extend this idea by feeding GNN predictions 
into external optimizers or discovery based search procedures, transforming TE into a 
learn,predict,optimize loop. This line of work primarily treats routing as a 
supervised regression problem, with GNNs learning a differentiable mapping from 
graph + routing configuration to performance.

A second stream couples GNN encoders with deep reinforcement learning (DRL) to 
directly learn routing or TE policies. G-Routing proposes a flexible online routing 
framework where a GNN encodes the current topology, traffic matrix, and routing state, 
predicting future path-level metrics that a DRL agent then uses to choose routes. \cite{wei2023grouting} 
Experiments show that G-Routing converges faster and yields lower jitter and more 
stable performance than shortest-path and traditional TE baselines under traffic 
dynamics and failures. \cite{wei2023grouting} Similarly, GROM (Generalized Routing Optimization Method) 
uses a GNN+DRL design to learn traffic-splitting policies; the authors encode split 
ratios as link features and discretize the continuous action space into classification, 
improving training stability and generalization to unseen topologies. \cite{Ding2024grom}


More recent works \cite{zhou2025trafficengineeringlargescalenetworks} explicitly target generalization across networks and demand patterns 
by designing GNN architectures aligned with classical TE formulations. 
A representative approach constructs a TE-GNN whose graph includes paths, 
source–destination demand pairs, links, and a global objective node, effectively 
embedding the linear program structure of TE into a heterogeneous graph. Message 
passing over this graph allows the GNN to learn correlations between paths, links, 
traffic and the global objective (e.g., maximum link utilization), and experiments 
indicate that such models can generalize TE policies to unseen topologies while 
approximately respecting constraints.\cite{zhou2025trafficengineeringlargescalenetworks} 
Jiang et al. (2024) \cite{jiang2024graph} broader survey situates these 
models within a larger family of GNN-based control and optimization methods, 
highlighting that TE and routing are a well studied application for GNNs.

\section{Taxonomy on Adversarial attacks on GNNs}

This subchapter will present contemporary research on adversarial attacks 
against Graph Neural Networks (GNNs), caracterizing them in 6 categories. 
The classification scheme addresses the unique challenges posed by 
graph-structured data and the message-passing mechanisms of GNNs that attackers
must address in their strategies.


\begin{figure}[h]
    \centering
    \includegraphics[width=0.8\textwidth]{Taxonomy_GNN.png}
    \caption{Diagram of Adversarial Attack characteristics on Graph Neural Networks}
    \label{fig:gnn_taxonomy}
\end{figure}



\subsection{Attacker knowledge}

Adversarial attacks on GNNs can be categorized based on the attacker's 
knowledge:

\begin{table}[htbp]
\centering
\small
\setlength{\tabcolsep}{4pt}
\caption{Attacker's Knowledge (Information Access Level)}
\label{tab:attacker_knowledge}
\begin{tabularx}{\textwidth}{XXX}
\toprule
\textbf{White-Box} & \textbf{Gray-Box} & \textbf{Black-Box} \\
\midrule
Complete model access: parameters, architecture, training data, labels & 
Limited access: training data and labels available; model parameters obscured & 
Minimal access: only adjacency matrix, feature matrix, and query access to model outputs \\
\addlinespace
Gradient-based optimization feasible; represents worst-case vulnerability assessment & 
Requires surrogate model training & 
Query-efficient methods: reinforcement learning, genetic algorithms, or transfer from surrogate models \\
\addlinespace
\bottomrule
\end{tabularx}
\end{table}

\newpage
\subsection{Attack Timing}
Adversarial attacks on GNNs can be categorized by when they are deployed:

\begin{itemize}
    \item \textbf{Training-Time Attacks (Poisoning Attacks)}: Adversarial examples are injected 
    into training data to manipulate the model's learned representations. \cite{luo2024survey,wei2020adversarial}
    \item \textbf{Test-Time Attacks (Evasion Attacks)}: Adversarial examples are introduced
    during test time to interfere with model predictions. \cite{luo2024survey,wei2020adversarial}
\end{itemize}

\subsection{Attack Strategy (Perturbation Goal)}

Adversarial attacks on GNNs can be categorized based on their perturbation goals:

\begin{itemize}
    \item \textbf{Topology Attacks}: Manipulate the graph structure by adding or removing edges to mislead the GNN.
    \item \textbf{Feature Attacks}: Alter node features to deceive the GNN's predictions.
    \item \textbf{Hybrid Attacks}: Combine topology and feature perturbations to maximize attack effectiveness.
\end{itemize}

\subsection{Attack Manipulation}

Adversarial attacks on GNNs can be categorized based on the manipulation method
whose effectiveness is generally evaluated by stealth rate while maintaining success rates:

\begin{itemize}
    \item \textbf{Addition}: Introduce new nodes or edges to the graph.
    \item \textbf{Removal}: Delete existing nodes or edges from the graph.
    \item \textbf{Rewiring}: A mix of addition and removal.
\end{itemize}

\textbf{Rewiring} operations tend to preserve graph statistics 
(node count, edge count, degree distribution) while enabling structural 
perturbations, making them significantly less detectable. \cite{11261632}

\subsection{Attack Specificity}

Adversarial attacks on GNNs can be categorized based on their specificity:

\begin{itemize}
    \item \textbf{Targeted Attacks}: Aim to degrade specific model functions or characteristics.
    \item \textbf{Untargeted Attacks}: Seek to degrade overall model performance without targeting specific functions.
\end{itemize}

\subsection{Task level}

Adversarial attacks on GNNs can be categorized based on the level in which they operate:

\begin{itemize}
    \item \textbf{Node-Level Attacks}: Target individual nodes to mislead node classification or regression tasks. A key vulnerability at this level in GNNs is the assumption of local homophily. \cite{luo2024survey}
    \item \textbf{Link-Level Attacks}: Focus on specific edges to disrupt link prediction tasks. At this level a common problem is the triadic closure disruption where nodes sometimes create links between them despite having another node which connects them, effectively creating a "triangle" \cite{11261632}
    \item \textbf{Graph-Level Attacks}: Aim to alter entire graph structures to affect graph classification tasks. An exploitable vulnerability at this level is the graph pooling instability. \cite{11261632}
\end{itemize}

\subsection{Examples and characteristics of adversarial attacks on GNNs}

This is by no means a thorough list of adversarial attacks on GNNs, but should give a good overview of the state of the art.

\begin{table}[htbp]
\centering
\scriptsize
\caption{Summary of Adversarial Attacks on GNNs}
\label{tab:attack_matrix}
\begin{tabularx}{\textwidth}{lccccccX}
\toprule
\textbf{Attack Method} & \textbf{Knowledge} & \textbf{Capability} & \textbf{Strategy} & \textbf{Manipulation} & \textbf{Goal} & \textbf{Task} \\
\midrule

Nettack \cite{Nettack} & White-box & Evasion & Topology & Rewiring & Targeted & Node \\

Metattack \cite{metaattackgnn} & White-box & Poisoning & Topology & Rewiring & Untargeted & Node \\

FGA \cite{chen2018fast} & White-box & Evasion & Topology & Rewiring & Targeted & Node \\

RL-S2V \cite{RL-S2V} & Black-box & Evasion & Topology & Rewiring (edge flips) & Targeted & Node/Link \\

PGD-Attack \cite{xu2019topology} & White-box & Evasion & Topology & Rewiring & Both & Node \\

NIPA \cite{NIPA} & Gray-box & Poisoning & Hybrid & Inject nodes & Targeted & Node \\

ReWatt \cite{ma2019attacking} & Black-box & Evasion & Topology & Rewiring & Untargeted & Graph \\

IG-FGSM \cite{wu2019adversarial}& White-box & Evasion & Hybrid & Rewiring(Feature+Edge) & Both & Node \\

\bottomrule
\end{tabularx}
\end{table} \cite{luo2024survey,wei2020adversarial}

In the literature on graph adversarial learning,
Nettack and Metattack are typically implemented with full access to model 
parameters and gradients, which designates them as white-box attacks under 
standard threat-model definitions \cite{Nettack,metaattackgnn}. However, some
more recent analyses decribe it as gray-box when 
emphasizing that the optimized perturbations are later 
transferred to a separate victim models. \cite{luo2024survey,wei2020adversarial}

\subsection{Transferability-Focused Attacks}

Recent work demonstrates that adversarial examples transfer across GNN 
architectures, datasets, and even task levels. GOttack leverages graph 
orbit learning to generate universal perturbations effective against 
multiple GNN variants. Migratability studies show 70-85\% attack success 
rate when transferring between GCN, GAT, and GraphSAGE. \cite{11004852,11199215}


